% !TEX program = xelatex
\documentclass[10pt,aspectratio=169]{beamer}

\usepackage[T1]{fontenc}
\usepackage{lmodern}

% Chinese (Simplified) support for the translated deck. xeCJK routes CJK
% codepoints to a CJK font while leaving Latin text/math on the fontspec
% faces set per-deck; this is lighter-touch than full `ctex`, which fights
% beamer/metropolis's own fontspec setup. Requires XeLaTeX (already the
% engine for these decks) and Noto CJK SC (installed system-wide).
\usepackage{xeCJK}
\setCJKmainfont{Noto Sans CJK SC}
\setCJKsansfont{Noto Sans CJK SC}
\setCJKmonofont{Noto Sans Mono CJK SC}

% Header bar matches the rafael-sharon toronto-may-2026 reference deck:
% miniframes navigation with subsection ticks, dark-teal section/subsection
% bars at the top of every content frame.
\definecolor{bg_subsec}{RGB}{43, 66, 70}
\definecolor{bg_sec}{RGB}{51, 77, 83}

\usetheme{metropolis}
\useoutertheme[subsection=true]{miniframes}
\setbeamercolor{subsection in head/foot}{fg=white, bg=bg_subsec}
\setbeamercolor{section in head/foot}{fg=white, bg=bg_sec}

\usepackage{appendixnumberbeamer}

\usepackage{amsmath, amssymb, amsthm}
\usepackage{mathtools}
\usepackage{bm}
\usepackage{graphicx}
\usepackage{booktabs}

% Code listings
\usepackage{listings}
\usepackage{xcolor}
\definecolor{jl-keyword}{HTML}{8959A8}
\definecolor{jl-string}{HTML}{718C00}
\definecolor{jl-comment}{HTML}{8E908C}
\lstdefinelanguage{Julia}{
  keywords={if,else,elseif,end,for,while,function,return,using,import,
            module,struct,abstract,type,export,begin,let,do,in,true,false,
            const,where},
  morecomment=[l]{\#},
  morestring=[b]",
  sensitive=true,
}
\lstset{
  basicstyle=\ttfamily\footnotesize,
  language=Julia,
  keywordstyle=\color{jl-keyword}\bfseries,
  stringstyle=\color{jl-string},
  commentstyle=\color{jl-comment},
  showstringspaces=false,
  breaklines=true,
  tabsize=2,
  captionpos=b,
  frame=single,
  framerule=0pt,
  backgroundcolor=\color{black!3},
  extendedchars=true,
  literate=
    {β}{{$\beta$}}1 {α}{{$\alpha$}}1 {γ}{{$\gamma$}}1 {δ}{{$\delta$}}1
    {ε}{{$\varepsilon$}}1 {ζ}{{$\zeta$}}1 {η}{{$\eta$}}1 {θ}{{$\theta$}}1
    {κ}{{$\kappa$}}1 {λ}{{$\lambda$}}1 {μ}{{$\mu$}}1 {ν}{{$\nu$}}1
    {ξ}{{$\xi$}}1 {π}{{$\pi$}}1 {ρ}{{$\rho$}}1 {σ}{{$\sigma$}}1
    {τ}{{$\tau$}}1 {φ}{{$\varphi$}}1 {χ}{{$\chi$}}1 {ψ}{{$\psi$}}1
    {ω}{{$\omega$}}1 {Λ}{{$\Lambda$}}1 {Δ}{{$\Delta$}}1 {Σ}{{$\Sigma$}}1,
}

% Common math macros (kept in sync with paper/main.tex)
\newcommand{\R}{\mathbb{R}}
\newcommand{\E}{\mathbb{E}}
\newcommand{\Prob}{\mathbb{P}}
\newcommand{\norm}[1]{\left\| #1 \right\|}
\newcommand{\Vstar}{V^{\star}}

\newcommand{\p}[1]{\left( #1 \right)}
\newcommand{\psq}[1]{\left[ #1 \right]}

% Section page: use the long-form section name (\insertsection) instead of
% the short-form (\insertsectionhead) so \section[short]{long} renders
% `long` on the section-title slide and `short` in the header. Cloned
% from metropolis's `progressbar` section page template.
\setbeamertemplate{section page}{%
  \centering%
  \begin{minipage}{22em}%
    \raggedright%
    \usebeamercolor[fg]{section title}%
    \usebeamerfont{section title}%
    \insertsection\\[-1ex]%
    \usebeamertemplate*{progress bar in section page}%
    \par%
    \ifx\insertsubsection\@empty\else%
      \usebeamercolor[fg]{subsection title}%
      \usebeamerfont{subsection title}%
      \insertsubsection%
    \fi%
  \end{minipage}%
  \par%
  \vspace{\baselineskip}%
}

% Project-specific macros
\newcommand{\Vstart}{V^{\mathrm{start}}}
\newcommand{\Vend}{V^{\mathrm{end}}}
\newcommand{\Vstarttil}{\widetilde{V}^{\mathrm{start}}}
\newcommand{\Vendtil}{\widetilde{V}^{\mathrm{end}}}
\newcommand{\Vstarthat}{\widehat{V}^{\mathrm{start}}}
\newcommand{\Vendhat}{\widehat{V}^{\mathrm{end}}}
\newcommand{\Lstart}{\Lambda^{\mathrm{start}}}
\newcommand{\Lend}{\Lambda^{\mathrm{end}}}
\newcommand{\xstart}{x^{\mathrm{start}}}
\newcommand{\xend}{x^{\mathrm{end}}}
\newcommand{\CV}{\mathcal{V}}
\newcommand{\CL}{\mathcal{L}}


\usepackage{fontspec}

% Override Metropolis's Fira default — use Computer Modern (Latin Modern under
% XeLaTeX) so the deck compiles deterministically without Fira installed.
\setmainfont{Latin Modern Roman}
\setsansfont{Latin Modern Sans}
% DejaVu Sans Mono (not Latin Modern Mono) for the monospace face because
% LM Mono lacks Greek glyphs — β/γ inside lstlisting and \texttt{} would be
% silently dropped otherwise.
\setmonofont{DejaVu Sans Mono}[Scale=0.85]

\title{第一讲：引论与动态建模}
\subtitle{异质性主体宏观经济学的计算方法}
\author{孙振越 (Jeffrey Sun)}
\institute{多伦多大学}
\date{2026 年 5 月 11 日}

\begin{document}

\maketitle

\begin{frame}{模型结构}
  本课程幻灯片：在 Julia 中求解含或不含总量冲击的异质性主体（HA）一般均衡（GE）模型。

  我们讨论具有如下形式的离散时间一般均衡模型：
  \begin{enumerate}
    \item 由若干\textbf{阶段}（Stages）构成的前瞻性家庭模块
    \item 一组描述经济其余部分（如厂商、政府等）的方程，通常表示为函数的\textbf{有向无环图}（DAG）。
    \item 一组\textbf{外层循环}，用于求解诸如均衡、稳态、校准、转移或全局解
  \end{enumerate}
\end{frame}

\begin{frame}{课程结构}
  \begin{itemize}
    \item 第 1--4 讲：阶段。
    \item 第 5--8 讲：外层循环。
    \begin{itemize}
      \item 最终引出“神经值函数迭代”（nVFI）全局求解算法
    \end{itemize}
  \end{itemize}
\end{frame}

% =============================================================
\section{动机}
% =============================================================

\begin{frame}{异质性主体：动机一}
  \begin{itemize}
    \item 经验上平均边际消费倾向（MPC）较高，而平均储蓄\textbf{也}较高
    \item 典型的代表性主体模型难以轻易解释这一点。
    \item 然而，若贫困家庭具有高 MPC、低储蓄，而富裕家庭具有低 MPC、高储蓄，则\textbf{异质性主体}模型便可加以解释。
  \end{itemize}
\end{frame}

\begin{frame}{异质性主体：动机二}
\begin{itemize}
  \item 政策（例如利率）的变化会对不同的人产生\textbf{不同}的影响
  \item 一项政策的\textbf{总量}效应取决于它针对的对象，因而各不相同。
  \item 我们关心政策在分布层面和总量层面的后果。
\end{itemize}

\end{frame}

% =============================================================
\section{家庭问题：两期}
% =============================================================

\begin{frame}{微观基础：单个家庭}
  \begin{itemize}
    \item 宏观经济学研究的是个体行为如何塑造总量结果。
    \item 要理解宏观结果的根源，就需要研究个体行为的来源。
  \end{itemize}
\end{frame}

\begin{frame}{两期模型}
  家庭在第 0 期以\textbf{财富} $b_0$ 开始。在第 1 期，家庭获得收入 $z_1$。
  \[
    \max_{c_0, c_1 \ge 0} \; u(c_0) + \beta u(c_1)
  \]
  \[
    \text{s.t.}\quad b_0 + \frac{z_1}{R} \geq c_0 + \frac{c_1}{R}.
  \]
  \begin{itemize}
    \item 若今天储蓄 \$1，明天便可花费 \$$R$，但效用被 $\beta$ 折现。
    \item 拉格朗日函数：$\mathcal{L} = u(c_0) + \beta u(c_1) - \lambda\p{c_0 + \frac{c_1}{R} - b_0 - \frac{z_1}{R}}$。
  \end{itemize}
\end{frame}

\begin{frame}{欧拉方程}
  拉格朗日一阶条件：
  \begin{align*}
    u'(c_0) &= \lambda, \\
    \beta u'(c_1) &= \frac{\lambda}{R}.
  \end{align*}
  整理可得，
  \[
    \boxed{\frac{u'(c_0)}{u'(c_1)} = \beta R.}
  \]
  在最优处，对于今天多消费一单位、还是把它储蓄以换取明天的 $R$ 单位，家庭是\textbf{无差异}的。
\end{frame}

\begin{frame}{示例：对数效用下的解析解}
  设 $u(c) = \log c$。则 $c_1 = \beta R c_0$，且
  \[
    c_0^\star = \frac{b_0 + \frac{z_1}{R}}{1+\beta},
    \qquad
    c_1^\star = \frac{\beta R \, b_0 + \frac{z_1}{R}}{1+\beta}.
  \]
  当 $\beta = 0.96$、$R = 1.04$、$b_0 = 10$、$z_1 = 2$ 时：
  \[
    c_0^\star \approx 6.08, \;\; c_1^\star \approx 6.07.
  \]
  这里存在\textbf{消费平滑}。收入是起伏不平的 $(10, 2)$，而消费则更为平滑，为 $(6.08, 6.07)$。
\end{frame}

\begin{frame}{策略函数}
  \textbf{策略函数}将最优的第 0 期消费表示为\textbf{初始财富} $b_0$ 的函数，
  \[
    c_0^\star(b_0) = \frac{b_0 + z_1/R}{1+\beta}.
  \]
  这样，我们便得到了一个关于家庭行为、以及\textbf{家庭行为}如何产生的模型。
\end{frame}

\begin{frame}{CRRA 偏好}
  一个标准的假设是常相对风险厌恶（CRRA）效用，
  \[
    u(c) = \begin{cases}
      \frac{c^{1-\gamma}-1}{1-\gamma}, & \gamma > 0, \gamma \neq 1, \\
      \log c, & \gamma = 1.
    \end{cases}
  \]
  效用关于 $c$ \textbf{递增}，
  \[
    \frac{\partial u(c)}{\partial c} = c^{-\gamma} > 0.
  \]
  效用是\textbf{凹}的：边际效用关于 $c$ \textbf{递减}，
  \[
    \frac{\partial^2 u(c)}{\partial c^2} = -\gamma c^{-\gamma-1} < 0.
  \]
\end{frame}

\section{家庭问题：$T$ 期}

\begin{frame}{$T$ 期模型}
  终身效用由下式给出，
  \[
    \max_{c_0, \ldots, c_{T-1}} \sum_{t=0}^{T-1} \beta^t u(c_t),
    \quad b_{t+1} = R(b_t - c_t) + z_{t+1}.
  \]
  求解拉格朗日函数可得 $T-1$ 个欧拉方程，
  \[
    u'(c_t) = \beta R \, u'(c_{t+1}).
  \]
  如果未来收入被完全预见，仍然可以相当容易地求解。

  但如果收入是随机的呢？
\end{frame}

\begin{frame}{随机收入}
  \begin{itemize}
    \item 假设 $\{z_t\}$ 服从转移矩阵为 $\Pi$ 的马尔可夫链
    \item 给定今天的 $z_t$，我们知道 $z_{t+1}$ 的\emph{条件分布}，但不知道其实现值
    \item 在每一期 $t$，家庭在不知道 $z_{t+1}$ 的情况下选择 $c_t$
  \end{itemize}
\end{frame}

\begin{frame}{随机欧拉方程}
  当 $T = 2$ 时，我们可以代入，
  \[
    c_1 = R(b_0 - c_0) + z_1.
  \]
  于是问题就是，
  \[
    \max_{c_0} u(c_0) + \beta \E\left[\, u\left(R(b_0 - c_0) + z_1\right) \right],
  \]
  关于 $c_0$ 的一阶条件为，
  \[
    \boxed{\;u'(c_0) = \beta R \, \E\left[ u'(c_1) \right].}
  \]
  现在欧拉方程中包含了对\emph{不确定}的 $c_1$ 取期望。
\end{frame}

\begin{frame}{无限期问题（$T = \infty$）}
  说来或许令人意外，如果我们取 $T\to\infty$，问题实际上反而稍微简单一些：
  \begin{align*}
    \max \E_0 & \ \sum_{t=0}^{\infty} \beta^t u(c_t)\\
    \text{s.t.}\quad b_{t+1} &= R(b_t - c_t) + z_{t+1}.
  \end{align*}
  但我们是在对什么取 $\max$ 呢？严格来说，我们需要找到一整个序列的\textbf{策略函数}，
  \[
    c_t = c_t(z_0, \ldots, z_t) \quad \forall t,
  \]
  但这看起来几乎是不可能的。所幸，我们可以通过重新表述问题来加以简化。

  （注：此处略去横截性条件的问题。）
\end{frame}

% =============================================================
\section{递归形式}
% =============================================================

\begin{frame}{Bellman 的第一个观察}
  在每一时刻 $t$，家庭只做出一个选择，然后等待，

  \[
    \max_{c_t} = u(c_t) + \beta \E_t\left[ \max \E_{t+1}\sum_{s=t+1}^\infty \beta^t u(c_s) \right]
  \]
  \[
    s.t. \quad b_{s+1} = R(b_s - c_s) + z_{s+1} \quad\forall s \ge t.
  \]

  重要的是，时刻 $t$ 的问题与时刻 $t+1$ 的问题看起来\textbf{几乎}完全相同。

  唯一可能的差异在于 $b_t$ 和 $z_t$。这些就是\textbf{状态变量}。
\end{frame}

\begin{frame}{状态变量}
  \begin{itemize}
    \item 非正式地说，\textbf{状态变量}就是那些使各个时期（以及各项决策）彼此之间产生实质性差异的变量。
    \item 在这里，状态变量是财富 $b_t$ 和收入 $z_t$。给定这些变量，家庭未来的条件期望效用和最优决策便被唯一确定。
  \end{itemize}
\end{frame}

\begin{frame}{Bellman 的第二和第三个观察}
  Bellman 还意识到另外两件事：
  \begin{enumerate}
    \item[2.] 可以将第一个洞见表述为一个关于\textbf{值函数} $V$ 的方程。
    \item[3.] 如果你知道这个\textbf{值函数}，就可以计算出\textbf{策略函数}，并（在数值上）求解整个模型。
  \end{enumerate}
  用计算机科学的术语来说，这是\textbf{缓存}的一个例子。
\end{frame}

\begin{frame}{值函数与贝尔曼方程}
  将\emph{值函数} $V(b, z)$ 定义为给定状态 $(b, z)$ 时未来效用的\textbf{净现值}。
  \begin{align*}
    V(b, z) &= \max_{c\in[0,b]} \mathbb{E}_0 \sum_{s=0}^\infty \beta^{s} u(c_s)\\
    \text{s.t.} \quad b_{s+1} &= R(b_s - c_s) + z_{s+1}\ \forall s.
  \end{align*}
  $\quad$（前提是所有未来的选择也都将以最优方式做出。）

  那么 $V$ 满足\textbf{贝尔曼方程}
  \[
    \boxed{V(b, z) = \max_{c \in [0, b]} \; u(c) + \beta \, \E\!\left[\, V(b', z') \,\big|\, z \,\right]}
  \]
  s.t. $b' = R(b-c) + z'$
\end{frame}

\begin{frame}{值函数就是你所需要的一切}
  如果我们知道 $V$，就可以计算出\textbf{策略函数}，并求解整个（家庭）模型，
  \[
  c^\star(b, z) = \arg\max_{c \in [0, b]} u(c) + \beta \, \E\!\left[\, V(R(b-c) + z,\, z') \mid z \,\right].
  \]
\end{frame}

\begin{frame}{我们如何找到 $V$？}
  \begin{itemize}
    \item 不过，我们仍然面临一个问题。我们要找的不只是一个数，而是要找到一整个\textbf{函数} $V$。
    \item 我们可以通过“值函数迭代”（VFI）来做到这一点。
    \item 为了理解这一点：先设想我们猜测了一个值函数，而且猜\textbf{错}了。
  \end{itemize}

\end{frame}

\begin{frame}{作用于任意 $v$ 的贝尔曼算子}
  假设我们猜测了一个\textbf{函数} $v$。即使 $v$ 可能是\textbf{错}的，它仍然会预测出某个策略，
  \[
    c^*(b, z; v) = \arg\max_{c \in [0, b]} u(c) + \beta \, \E\!\left[\, v(R(b-c) + z',\, z') \mid z \,\right].
  \]
  如果我们用 $v$ 来\textbf{向前看}，便可得到一个\textbf{不同的函数} $\mathcal{T}v$，
  \[
    (\mathcal{T} v)(b, z) \;:=\; \max_{c \in [0, b]} \; u(c) + \beta \, \E\!\left[\, v(R(b-c) + z',\, z') \,\big|\, z \,\right].
  \]
  $v$ 在某种意义上代表了对处于各种状态 $(b, z)$ 之价值的\textbf{信念}。

  $\mathcal{T}v$ 则代表了你通过向前看一期来推理后得到的\textbf{更新后的信念}。
\end{frame}

\begin{frame}{贝尔曼算子}
  \begin{itemize}
    \item 把 $\mathcal{T}$ 看作一个把\textbf{函数}映射为\textbf{函数}的\textbf{算子}：$v \mapsto \mathcal{T}v$。
    \item 这就是\textbf{贝尔曼算子}。
    \item 贝尔曼方程可以改写为，
    \[
      \boxed{\; V \;=\; \mathcal{T} V. \;}
    \]
    \item 我们可以通过寻找 $\mathcal{T}$ 的\textbf{不动点}来求解家庭问题。
  \end{itemize}
\end{frame}

\begin{frame}{存在性与唯一性}
  我们不会证明这一点，但值得指出的是，巴拿赫不动点定理意味着，在温和的条件下，$\mathcal{T}$ 存在\textbf{唯一}的不动点 $V$。

  其想法在于 $\mathcal{T}$ 是一个压缩映射：
  \[
    \norm{\mathcal{T} v_1 - \mathcal{T} v_2}_\infty
      \le \beta \norm{v_1 - v_2}_\infty.
  \]
  巴拿赫定理意味着：
  \begin{enumerate}
    \item 存在唯一的不动点 $V$。
    \item 从任意初始 $v_0$ 出发迭代 $v_{k+1} = \mathcal{T} v_k$，都将以 $\beta$ 的速率收敛到 $V$。
  \end{enumerate}
\end{frame}

\begin{frame}{再看一遍：贝尔曼方程}
\[
  \boxed{\;
    V(\text{状态}) = \max_{\text{行动}}\ \text{收益}(\text{状态, 行动}) + \beta \, \E\!\left[\, V(\,\text{下一状态}\,) \,\big|\, \text{状态, 行动} \,\right]
  \;}
\]
本课程余下的内容基本上都是这个方程的各种变体。
\end{frame}

% =============================================================
\section{后续讲次}
% =============================================================

\begin{frame}{后续讲次}
\centering
\footnotesize
\setlength{\tabcolsep}{4pt}
\renewcommand{\arraystretch}{0.9}
\begin{tabular}{c}
VFI 的理论 \quad (L1) \\
$\downarrow$ \\
Julia 中的 VFI \quad (L2) \\
$\downarrow$ \\
异质性 \quad (L3) \\
$\downarrow$ \\
一般均衡 \quad (L4) \\
$\downarrow$ \\
转移动态 \quad (L5) \\
$\downarrow$ \\
空间与校准 \quad (L6) \\
$\downarrow$ \\
总量不确定性 \quad (L7) \\
$\downarrow$ \\
\alert{神经值函数迭代} \quad (L8)
\end{tabular}
\end{frame}

\begin{frame}{Notebook}
  每一讲都配有一个对应的 Jupyter notebook。
  \begin{enumerate}
    \item Julia 1.10+\\
      \footnotesize\url{https://julialang.org/downloads/}
    \item \normalsize VS Code 加 Julia 扩展。
    \item \texttt{using Pkg; Pkg.add("IJulia")}。
    \item Jupyter（Anaconda 或 \texttt{pip install jupyter}）。
  \end{enumerate}
\end{frame}

\end{document}
